% GENERATED FILE. Edit canonical lessons, then run npm run generate:books.
% canonical-source-hash: fnv1a64:5dcfaba7820092b7
% canonical-lessons: SA-C84-temple, SA-C84-hospital, SA-C84-field, SA-C84-farmer, SA-C84-merchant

\chapter{Temple, Hospital and Field}
\label{ch:sa-temple-hospital-field}
\hlchaptermodality{84}

\begin{chapteropening}
\textbf{By the end of this chapter:} \emph{I can name a temple, a hospital and a field, a farmer and a merchant.}

Complete the last lesson of chapter 84: i can name a temple, a hospital and a field, a farmer and a merchant.
\end{chapteropening}

\section[mandiram]{\sk{मन्दिरम्} (mandiram) — a temple}

\label{lesson:SA-C84-temple}



\begin{quote}
Before the new one: say the Sanskrit for a lesson, then the Sanskrit for an exam.
\end{quote}

\subsection*{You'll want to know: \sk{मन्दिरम्}}
\textbf{\sk{मन्दिरम्}} (\emph{mandiram}) — \textquotedblleft{}a temple\textquotedblright{}. \textbf{\sk{मन्दिरम्} \sk{कुत्र}?} — \textquotedblleft{}where is the temple?\textquotedblright{}

\begin{grammarlens}[title={What you've built}]
A place of worship.
\end{grammarlens}

\subsection*{Your turn}
\begin{itemize}
\raggedright
  \item \emph{Say it:} \emph{mandiram}
  \item \emph{Say it:} \emph{mandiram}, once more
  \item \emph{Say it:} ask \emph{mandiram kutra?}
  \item \emph{Recall:} say the Sanskrit for a lesson, then the Sanskrit for an exam, then say \emph{mandiram} again
\end{itemize}

\begin{tcolorbox}[breakable,colback=teal!4,colframe=teal!35!black,title={Before you move on}]
What does \sk{मन्दिरम्} mean? (\textquotedblleft{}A temple\textquotedblright{}.) Say it once more.
\end{tcolorbox}
\section[cikitsālayaḥ]{\sk{चिकित्सालयः} (cikitsālayaḥ) — a hospital}

\label{lesson:SA-C84-hospital}



\begin{quote}
Before the new one: say the Sanskrit for an exam, then the Sanskrit for a temple.
\end{quote}

\subsection*{You'll want to know: \sk{चिकित्सालयः}}
\textbf{\sk{चिकित्सालयः}} (\emph{cikitsālayaḥ}) — \textquotedblleft{}a hospital\textquotedblright{}: \textbf{\sk{चिकित्सा}}, \textquotedblleft{}treatment\textquotedblright{}, with the same \textbf{\sk{आलयः}}, \textquotedblleft{}home\textquotedblright{}, as \textbf{\sk{विद्यालयः}}.

\begin{grammarlens}[title={What you've built}]
Treatment plus home.
\end{grammarlens}

\subsection*{Your turn}
\begin{itemize}
\raggedright
  \item \emph{Say it:} \emph{cikitsālayaḥ}
  \item \emph{Say it:} \emph{cikitsālayaḥ}, once more
  \item \emph{Say it:} ask \emph{cikitsālayaḥ kutra?}
  \item \emph{Recall:} say the Sanskrit for an exam, then the Sanskrit for a temple, then say \emph{cikitsālayaḥ} again
\end{itemize}

\begin{tcolorbox}[breakable,colback=teal!4,colframe=teal!35!black,title={Before you move on}]
What does \sk{चिकित्सालयः} mean? (\textquotedblleft{}A hospital\textquotedblright{}.) Say it once more.
\end{tcolorbox}
\section[kṣetram]{\sk{क्षेत्रम्} (kṣetram) — a field}

\label{lesson:SA-C84-field}



\begin{quote}
Before the new one: say the Sanskrit for a temple, then the Sanskrit for a hospital.
\end{quote}

\subsection*{You'll want to know: \sk{क्षेत्रम्}}
\textbf{\sk{क्षेत्रम्}} (\emph{kṣetram}) — \textquotedblleft{}a field\textquotedblright{} of a \textbf{\sk{ग्रामः}}.

\begin{grammarlens}[title={What you've built}]
Where the rice grows.
\end{grammarlens}

\subsection*{Your turn}
\begin{itemize}
\raggedright
  \item \emph{Say it:} \emph{kṣetram}
  \item \emph{Say it:} \emph{kṣetram}, once more
  \item \emph{Say it:} say \emph{kṣetram}
  \item \emph{Recall:} say the Sanskrit for a temple, then the Sanskrit for a hospital, then say \emph{kṣetram} again
\end{itemize}

\begin{tcolorbox}[breakable,colback=teal!4,colframe=teal!35!black,title={Before you move on}]
What does \sk{क्षेत्रम्} mean? (\textquotedblleft{}A field\textquotedblright{}.) Say it once more.
\end{tcolorbox}
\section[kṛṣakaḥ]{\sk{कृषकः} (kṛṣakaḥ) — a farmer}

\label{lesson:SA-C84-farmer}



\begin{quote}
Before the new one: say the Sanskrit for a hospital, then the Sanskrit for a field.
\end{quote}

\subsection*{You'll want to know: \sk{कृषकः}}
\textbf{\sk{कृषकः}} (\emph{kṛṣakaḥ}) — \textquotedblleft{}a farmer\textquotedblright{}, who works the \textbf{\sk{क्षेत्रम्}}.

\begin{grammarlens}[title={What you've built}]
Who works the field.
\end{grammarlens}

\subsection*{Your turn}
\begin{itemize}
\raggedright
  \item \emph{Say it:} \emph{kṛṣakaḥ}
  \item \emph{Say it:} \emph{kṛṣakaḥ}, once more
  \item \emph{Say it:} say \emph{kṛṣakaḥ}
  \item \emph{Recall:} say the Sanskrit for a hospital, then the Sanskrit for a field, then say \emph{kṛṣakaḥ} again
\end{itemize}

\begin{tcolorbox}[breakable,colback=teal!4,colframe=teal!35!black,title={Before you move on}]
What does \sk{कृषकः} mean? (\textquotedblleft{}A farmer\textquotedblright{}.) Say it once more.
\end{tcolorbox}
\section[vaṇik]{\sk{वणिक्} (vaṇik) — a merchant}

\label{lesson:SA-C84-merchant}



\begin{quote}
Before the new one: say the Sanskrit for a field, then the Sanskrit for a farmer.
\end{quote}

\subsection*{You'll want to know: \sk{वणिक्}}
\textbf{\sk{वणिक्}} (\emph{vaṇik}) — \textquotedblleft{}a merchant\textquotedblright{}, who keeps the \textbf{\sk{आपणः}}.

\begin{grammarlens}[title={What you've built}]
Who keeps the shop.
\end{grammarlens}

\subsection*{Your turn}
\begin{itemize}
\raggedright
  \item \emph{Say it:} \emph{vaṇik}
  \item \emph{Say it:} \emph{vaṇik}, once more
  \item \emph{Say it:} say \emph{vaṇik}
  \item \emph{Recall:} say the Sanskrit for a field, then the Sanskrit for a farmer, then say \emph{vaṇik} again
\end{itemize}

\begin{tcolorbox}[breakable,colback=teal!4,colframe=teal!35!black,title={Before you move on}]
What does \sk{वणिक्} mean? (\textquotedblleft{}A merchant\textquotedblright{}.) Say it once more.
\end{tcolorbox}
